% Generated draft: requires mathematical and source review.
\documentclass[11pt]{article}
\usepackage[margin=25mm]{geometry}
\usepackage{fontspec}
\setmainfont{Latin Modern Roman}
\usepackage{amsmath,amssymb,mathtools,unicode-math}
\setmathfont{Latin Modern Math}
\usepackage{enumitem,needspace,xurl,hyperref}
\setlength{\parindent}{0pt}
\setlength{\parskip}{5pt}
\setlength{\emergencystretch}{3em}
\newcommand{\N}{\mathbb{N}}
\newcommand{\Z}{\mathbb{Z}}
\newcommand{\Q}{\mathbb{Q}}
\newcommand{\R}{\mathbb{R}}
\newcommand{\C}{\mathbb{C}}
\newcommand{\F}{\mathbb{F}}
\newcommand{\A}{\mathbb{A}}
\newcommand{\PP}{\mathbb{P}}
\newcommand{\E}{\mathbb{E}}
\newcommand{\eps}{\varepsilon}
\newcommand{\dd}{\,\mathrm{d}}
\begin{document}
\section*{UnsolvedMath ID: 30000369}
\textbf{Source problem code:} catalogue entry
\subsection{Definitions and mathematical statement}
Let $p:N\to M$ be a smooth surjection of manifolds, and let $A$ on $M$ and $B$ on $N$ be diffusion generators satisfying
\[
 B(f\circ p)=(Af)\circ p\qquad\text{for every }f\in C^\infty(M;\R).
\]
With $p^*f=f\circ p$, this is the differential-operator intertwining identity $Bp^*=p^*A$. The source asks to determine structural and spectral relationships between $A$ and $B$, including whether $B$ has a natural decomposition.

In the usual differential interpretation a diffusion generator is a second-order operator with nonnegative principal symbol and no killing term, so $A1=B1=0$. The source does not specify nondegeneracy, constant rank, completeness, measures, or boundary conventions; this standard interpretation must be confirmed if more general generators are intended.

A possible meaning of the decomposition request would be a canonical split into a part lifting $A$ and a part acting vertically, where a vertical operator $V$ annihilates pullbacks: $V(f\circ p)=0$ for all smooth $f$. If one writes $B=H+V$, the lifting condition would be $H(f\circ p)=(Af)\circ p$. This expresses a candidate type of decomposition, not a source-supplied theorem or an assertion of existence, uniqueness, positivity, commutation, or independence of the parts. The source supplies no horizontal distribution or preferred connection. A smooth surjection need not be a submersion, so no regular-fiber geometry is assumed without further hypotheses.

A spectral question requires closed realizations of $A$ and $B$ on specified function spaces, with domains, boundary conditions, and any invariant or reference measures. The formal pullback identity alone does not identify those realizations or specify whether the desired relation concerns eigenvalues, multiplicities, full spectra, spectral gaps, or semigroups. No particular spectral inclusion or equality is substituted for the report's broad request. Thus the explicit compatibility hypothesis is retained, while the intended structural result and spectral comparison remain to be recovered from the source.
\subsection{Short English statement}
The given diffusion on $N$ projects to the one on $M$ through the displayed pullback identity. The question asks for further structure, a natural decomposition, and spectral comparisons, but does not define the decomposition or the operator realizations needed for spectra.
\subsection{Sources}
\begin{itemize}
\item Local UnsolvedMath registry, permanent ID 30000369, code catalogue entry, ``Structural Relations Between Projected Diffusion Generators.'' Proposed year 2005; proposer unspecified.
\item \emph{Heat Kernels, Stochastic Processes and Functional Inequalities} (2006), Oberwolfach Reports, pp.\ 3061--3120, DOI 10.4171/owr/2005/54; original report \url{https://doi.org/10.4171/owr/2005/54}. The displayed report year 2006 differs from the DOI and proposed year 2005.
\item Statement status corrected\_verified; the supplied verification says raw maps, operators, and composition notation were enclosed in balanced math delimiters, retaining the task.
\item Literature triage checked 2026-08-21: open. It says the missing specifications of $A$, $B$, and the proposed relation leave an unresolved formulation rather than a verifiable theorem. The restored statement supplies the projection identity but not the additional target. Created 2026-08-25; updated 2026-09-16. Status and cited claims require independent verification.
\item Registry curation by Ulam AI / UnsolvedMath Contributors, CC BY 4.0; the underlying report retains its supplied provenance. Unreviewed statement-only draft.
\end{itemize}
\end{document}
